% !TEX root = web.tex

% Web variant of the blueprint, used by `leanblueprint` to produce the
% online HTML blueprint via plastex. Loads the same content but with the
% `blueprint` package which provides clickable Lean/dependency cross-links.

\documentclass{report}

\usepackage[utf8]{inputenc}
\usepackage{amsmath, amssymb, amsthm}
\usepackage{hyperref}
\usepackage[capitalize]{cleveref}

% Loads the real macros for \lean, \leanok, \uses, \notready, \proves
\usepackage{blueprint}

\newtheorem{theorem}{Theorem}[chapter]
\newtheorem{proposition}[theorem]{Proposition}
\newtheorem{lemma}[theorem]{Lemma}
\newtheorem{corollary}[theorem]{Corollary}
\newtheorem{remark}[theorem]{Remark}
\theoremstyle{definition}
\newtheorem{definition}[theorem]{Definition}
\newtheorem{example}[theorem]{Example}

% plasTeX renders thebibliography as a plain list with no heading, which leaves
% the references glued to the end of the last section. A part of their own
% gives them a heading, their own page (split-level 2) and an entry in the table
% of contents next to the six parts, outside the last one. It carries no \label:
% plasTeX binds a label to the last *numbered* node (Context.label via
% refstepcounter), so a label here would attach to the final theorem instead.
% The page is therefore named sectNNNN.html rather than after a label.
% The PDF build defines this differently -- see print.tex.
\newcommand{\bpbibheading}{\part*{References}}

% The introduction is an unnumbered chapter before the first part, with its own
% page and an entry in the table of contents. The PDF build defines this
% differently -- see print.tex.
\newcommand{\bpintroheading}{\chapter*{Introduction}}

% Parts carry roman numbers, as in the PDF, so that they stay apart from the
% arabic chapter numbers in the navigation.
\renewcommand{\thepart}{\Roman{part}}

% --- notation -------------------------------------------------------------
\newcommand{\Besov}[3]{B^{#1}_{#2,#3}}
\newcommand{\TL}[3]{F^{#1}_{#2,#3}}
\newcommand{\Sob}[2]{W^{#1}_{#2}}
\newcommand{\T}{\mathbb{T}}
\newcommand{\R}{\mathbb{R}}
\newcommand{\N}{\mathbb{N}}
\newcommand{\Z}{\mathbb{Z}}
\newcommand{\C}{\mathbb{C}}
\newcommand{\dist}{\operatorname{dist}}
\newcommand{\spn}{\operatorname{span}}

\title{Blueprint: Characterizations of Function Spaces}
\author{Mario Ullrich}

\begin{document}
\maketitle
% Links rendered below the title: the project home (the GitHub Pages site this
% blueprint is deployed to) and the source repository.
\home{https://mario-ullrich.github.io/Lean-FunctionSpaces/}
\github{https://github.com/mario-ullrich/Lean-FunctionSpaces}

% !TEX root = print.tex

\chapter{Approximation schemes}
\label{ch:schemes}

This blueprint follows the formalisation of approximation spaces in the sense of
Pietsch~\cite{Pietsch}, of Besov and Sobolev spaces defined through them, and of the
embedding theorems between these spaces. Each statement names the Lean declaration that
formalises it; a statement marked as formalised compiles without \texttt{sorry}.

Throughout, $X$ and $Y$ are real normed spaces. Complex spaces, such as $L_p$ spaces of
complex-valued functions, are regarded as real ones.

\begin{definition}[Approximation scheme]
\label{def:scheme}
\lean{ApproximationScheme}
\leanok
An \emph{approximation scheme} in $X$ is a sequence of subsets
$A_0 \subseteq A_1 \subseteq A_2 \subseteq \cdots$ of $X$ such that
\begin{enumerate}
  \item $A_0 = \{0\}$,
  \item $c\,a \in A_n$ for all $c \in \R$, $n \in \N$ and $a \in A_n$,
  \item $a + b \in A_{m+n}$ for all $m, n \in \N$, $a \in A_m$ and $b \in A_n$.
\end{enumerate}
\end{definition}

The elements of $A_n$ are the approximants of complexity $n$. The sets need not be linear
subspaces: in the $n$-term scheme of \Cref{def:nterm} they are unions of subspaces. A family
indexed by a degree enters with a strict inequality, $A_n = \{\text{degree} < n\}$, so that
$A_0 = \{0\}$ holds.

\begin{lemma}[Elementary properties]
\label{lem:scheme-basic}
\lean{ApproximationScheme.zero_mem, ApproximationScheme.nonempty,
  ApproximationScheme.neg_mem, ApproximationScheme.sub_mem, ApproximationScheme.sum_mem,
  ApproximationScheme.smul_set_eq, ApproximationScheme.sub_mem_pow_add_two}
\leanok
\uses{def:scheme}
Let $A$ be an approximation scheme in $X$. Then for all $m, n, k \in \N$:
\begin{enumerate}
  \item $0 \in A_n$; in particular $A_n$ is nonempty;
  \item $a - b \in A_{m+n}$ for all $a \in A_m$ and $b \in A_n$;
  \item $\sum_{i \in I} a_i \in A_{\sum_{i \in I} n_i}$ for every finite family with
    $a_i \in A_{n_i}$;
  \item $c\,A_n = A_n$ for every $c \in \R \setminus \{0\}$;
  \item $a - b \in A_{2^{k+2}}$ for all $a \in A_{2^{k+1}}$ and $b \in A_{2^k}$.
\end{enumerate}
\end{lemma}

\begin{proof}
\leanok
The first item follows from $A_0 = \{0\} \subseteq A_n$, the second from $-b = (-1)\,b$, and
the third by induction over $I$. For the fourth, $c^{-1} a \in A_n$ for $a \in A_n$. The last
follows from the second and $2^{k+1} + 2^k \leq 2^{k+2}$.
\end{proof}

\begin{definition}[Preimage and image of a scheme]
\label{def:comap}
\lean{ApproximationScheme.comap, ApproximationScheme.map}
\leanok
\uses{def:scheme}
Let $A$ be an approximation scheme in $X$. For an injective linear map $\iota : Y \to X$, the
sets $\iota^{-1}(A_n)$ form an approximation scheme in $Y$. For a linear map $T : X \to Y$, the
sets $T(A_n)$ form an approximation scheme in $Y$.
\end{definition}

Injectivity of $\iota$ is what gives $\iota^{-1}(\{0\}) = \{0\}$.

\section{The error of best approximation}

\begin{definition}[Error of best approximation]
\label{def:error}
\lean{ApproximationScheme.approxError}
\leanok
\uses{def:scheme}
Let $A$ be an approximation scheme in $X$. For $n \in \N$ and $f \in X$,
\[
  E_n(f) = \dist(f, A_n) = \inf_{a \in A_n} \|f - a\|.
\]
\end{definition}

\begin{lemma}[Calculus of the error]
\label{lem:error-basic}
\lean{ApproximationScheme.approxError_zero_left, ApproximationScheme.approxError_le_norm,
  ApproximationScheme.approxError_anti, ApproximationScheme.approxError_smul,
  ApproximationScheme.approxError_add_le, ApproximationScheme.abs_approxError_sub_le}
\leanok
\uses{def:error, lem:scheme-basic}
Let $A$ be an approximation scheme in $X$. For all $f, g \in X$, $c \in \R$ and
$m, n \in \N$:
\begin{enumerate}
  \item $E_0(f) = \|f\|$ and $E_n(f) \leq \|f\|$;
  \item $E_n(f) \leq E_m(f)$ whenever $m \leq n$;
  \item $E_n(c f) = |c| \, E_n(f)$;
  \item $E_{m+n}(f + g) \leq E_m(f) + E_n(g)$;
  \item $|E_n(f) - E_n(g)| \leq \|f - g\|$.
\end{enumerate}
\end{lemma}

\begin{proof}
\leanok
The first item holds because $A_0 = \{0\}$ and $0 \in A_n$, the second because the sets
increase, the third because $c A_n = A_n$ for $c \neq 0$. For the fourth, given
$\varepsilon > 0$ choose $a \in A_m$ and $b \in A_n$ with $\|f - a\| \leq E_m(f) +
\varepsilon$ and $\|g - b\| \leq E_n(g) + \varepsilon$; then $a + b \in A_{m+n}$ and the
triangle inequality gives the bound up to $2\varepsilon$. The last item is the
$1$-Lipschitz property of the distance to a set.
\end{proof}

A best approximant, an $a \in A_n$ with $\|f - a\| = E_n(f)$, need not exist, since the sets
$A_n$ need not be closed. The proofs work with near-best approximants as in the proof of the
fourth item.

\begin{lemma}[Convergence of the errors]
\label{lem:error-tendsto}
\lean{ApproximationScheme.tendsto_approxError_zero_iff}
\leanok
\uses{lem:error-basic}
Let $A$ be an approximation scheme in $X$ and $f \in X$. Then $E_n(f) \to 0$ as
$n \to \infty$ if and only if $f$ lies in the closure of $\bigcup_n A_n$.
\end{lemma}

\begin{proof}
\leanok
If $\|f - a\| < \varepsilon$ for some $a \in A_m$, then $E_n(f) < \varepsilon$ for all
$n \geq m$ since the errors decrease. Conversely $E_n(f) < \varepsilon$ gives such an $a$.
\end{proof}

\begin{lemma}[Functoriality of the error]
\label{lem:error-map}
\lean{ApproximationScheme.approxError_map_le}
\leanok
\uses{def:error}
Let $A$ and $B$ be approximation schemes in $X$ and $Y$, and let $T : X \to Y$ be a bounded
linear map with $T(A_n) \subseteq B_n$ for all $n$. Then for all $n \in \N$ and $f \in X$,
\[
  E^B_n(T f) \leq \|T\| \, E^A_n(f).
\]
\end{lemma}

\begin{proof}
\leanok
For $a \in A_n$, $E^B_n(Tf) \leq \|Tf - Ta\| \leq \|T\| \, \|f - a\|$; take the infimum
over $a$.
\end{proof}

\section{Two schemes}

\begin{definition}[Linear approximation]
\label{def:ofspan}
\lean{ApproximationScheme.ofSpan}
\leanok
\uses{def:scheme}
For a sequence $(e_i)_{i \in \N}$ in $X$, the sets $A_n = \spn\{e_i : i < n\}$ form an
approximation scheme.
\end{definition}

\begin{lemma}
\label{lem:ofspan}
\lean{ApproximationScheme.isClosed_ofSpan, ApproximationScheme.iUnion_ofSpan,
  ApproximationScheme.tendsto_approxError_ofSpan_iff}
\leanok
\uses{def:ofspan, lem:error-tendsto}
For the scheme of \Cref{def:ofspan}, every $A_n$ is closed, $\bigcup_n A_n = \spn\{e_i : i
\in \N\}$, and $E_n(f) \to 0$ if and only if $f$ lies in the closed span of $(e_i)$.
\end{lemma}

\begin{proof}
\leanok
The sets are finite-dimensional subspaces and therefore closed. The union of an increasing
family of spans is the span of the union. The last statement is \Cref{lem:error-tendsto}.
\end{proof}

\begin{definition}[$n$-term approximation]
\label{def:nterm}
\lean{ApproximationScheme.nTerm}
\leanok
\uses{def:scheme}
For a set $D \subseteq X$, the \emph{dictionary}, let $A_n$ be the set of linear combinations
of at most $n$ elements of $D$. These sets form an approximation scheme.
\end{definition}

The third axiom holds because a combination of $m$ and one of $n$ elements of $D$ together
use at most $m + n$ elements.

\begin{lemma}
\label{lem:nterm-ofspan}
\lean{ApproximationScheme.ofSpan_subset_nTerm, ApproximationScheme.approxError_nTerm_le_ofSpan}
\leanok
\uses{def:ofspan, def:nterm, lem:error-basic}
For a sequence $(e_i)$ in $X$, the linear scheme of $(e_i)$ is contained in the $n$-term
scheme of the dictionary $\{e_i : i \in \N\}$, and hence the $n$-term errors are at most the
linear ones.
\end{lemma}

\begin{proof}
\leanok
An element of $\spn\{e_i : i < n\}$ uses the $n$ elements $e_0, \dots, e_{n-1}$.
\end{proof}

\bpbibheading
\begin{thebibliography}{99}

\bibitem{BerghLofstrom}
J.~Bergh, J.~L\"ofstr\"om,
\emph{Interpolation Spaces. An Introduction},
Grundlehren der mathematischen Wissenschaften \textbf{223}, Springer-Verlag, Berlin, 1976,
\href{https://doi.org/10.1007/978-3-642-66451-9}{doi}.

\bibitem{DeVoreLorentz}
R.~A.~DeVore, G.~G.~Lorentz,
\emph{Constructive Approximation},
Grundlehren der mathematischen Wissenschaften \textbf{303}, Springer-Verlag, Berlin, 1993,
\href{https://doi.org/10.1007/978-3-662-02888-9}{doi}.

\bibitem{Nikolskii}
S.~M.~Nikol'ski\u{\i},
\emph{Approximation of Functions of Several Variables and Imbedding Theorems},
Grundlehren der mathematischen Wissenschaften \textbf{205}, Springer-Verlag, Berlin, 1975,
\href{https://doi.org/10.1007/978-3-642-65711-5}{doi}.

\bibitem{Pietsch}
A.~Pietsch,
\emph{Approximation spaces},
J.\ Approx.\ Theory \textbf{32} (1981), no.~2, 115--134,
\href{https://doi.org/10.1016/0021-9045(81)90109-X}{doi}.

\bibitem{Temlyakov}
V.~Temlyakov,
\emph{Multivariate Approximation},
Cambridge Monographs on Applied and Computational Mathematics \textbf{32}, Cambridge
University Press, Cambridge, 2018,
\href{https://doi.org/10.1017/9781108689687}{doi}.

\bibitem{Triebel}
H.~Triebel,
\emph{Theory of Function Spaces},
Monographs in Mathematics \textbf{78}, Birkh\"auser, Basel, 1983,
\href{https://doi.org/10.1007/978-3-0346-0416-1}{doi}.

\end{thebibliography}


\end{document}
